\chapter{De la prédiction au portefeuille}\label{ch:ml:from-prediction-to-portfolio}

La prévision de plus faible erreur de l'expérience du chapitre produit un portefeuille au ratio de Sharpe net de 0.19~; une prévision
à erreur légèrement supérieure en produit un de 1.52. La première a été ajustée sur tous les titres et a appris le signal fort
des titres que le desk ne peut pas négocier~; la seconde a été ajustée sur ceux qu'il peut négocier, où le signal est plus faible et
différent. Une prévision est une donnée d'entrée d'une décision, et la perte servant à l'entraîner doit tenir compte de ce que la décision
fait d'elle. Ce chapitre compare les trois façons de l'y amener~: entraîner la prévision seulement là où le portefeuille l'utilise, entraîner
la prévision à travers la règle de portefeuille (\omterm{def:ml:from-prediction-to-portfolio:dfl}{apprentissage orienté décision}), ou se passer de la prévision et apprendre directement la règle de portefeuille
(\omterm{def:ml:from-prediction-to-portfolio:ppp}{politiques de portefeuille paramétriques}). Il se conclut sur le prix des deux dernières~: elles ajustent un bruit que les moindres
carrés ignoreraient.

\section{Pourquoi la perte de la prévision n'est pas la perte de trading}

\begin{definition}[Prédire puis optimiser]\label{def:ml:from-prediction-to-portfolio:pto}
\emph{Prédire puis optimiser}\index{prédire puis optimiser} ajuste une prévision par une perte statistique (erreur quadratique moyenne, une
vraisemblance) et la transmet à un optimiseur qui choisit le portefeuille~; les deux étapes sont ajustées séparément.
\end{definition}

Le livre~7 (chapitre~15) mesurait la perte entre prévision et portefeuille par le coefficient de transfert~; ici, la perte se situe
plus en amont, dans la prévision elle-même. L'erreur quadratique moyenne pondère de la même manière tous les titres et toutes les périodes~; le profit du portefeuille
dépend de la prévision seulement là où l'optimiseur négocie, proportionnellement aux positions qu'il prend, et après coûts. Lorsque le
modèle de prévision est correctement spécifié, les moindres carrés estiment la vraie relation et l'écart n'a pas d'importance. Quand il
ne l'est pas, le modèle consacre sa flexibilité limitée aux erreurs qui comptent le plus pour la perte, qui ne sont pas forcément
celles qui comptent pour le portefeuille (Elmachtoub et Grigas, 2022).

Le panel du chapitre rend cela concret. Deux cents titres portent quatre caractéristiques qui suivent des processus AR(1) lents. Dans
les cent titres que le desk ne peut pas négocier (trop petits, soumis à restriction), la première caractéristique prédit
le rendement de la période suivante (0.3~\% par unité)~; dans les cent qu'il peut négocier, seule la seconde le fait, faiblement (0.08~\%). Les rendements comportent un bruit de 3~\% par période.
La prévision est linéaire par rapport aux caractéristiques, avec un coefficient unique pour tous les titres, si bien qu'elle ne peut pas être juste partout. La
règle de portefeuille est une optimisation moyenne-variance sur les titres négociables, avec une approximation quadratique du coût,
\[
w_{i,t} = \frac{\hat\mu_{i,t} + 2\kappa w_{i,t-1}}{\gamma\sigma^2 + 2\kappa},
\]
avec $\gamma = 20$, $\sigma = 3~\%$ et $\kappa$ égal au coût linéaire de 5 points de base que paient les rendements nets à chaque
variation de poids. Cinq paires de panels indépendants, de 1\,500 périodes chacune, fournissent les données d'entraînement et de test.

Les moindres carrés sur l'ensemble des titres attribuent un coefficient de 1.51 (en unités de 0.1~\%) à la première caractéristique et 0.43 à la seconde
(la \cref{fig:ml:e2e:coefs})~: il a surtout appris le signal non négociable. Hors échantillon, son erreur quadratique moyenne est inférieure à
celle des moindres carrés sur les seuls titres négociables (9.007 contre 9.031, en unités de $10^{-4}$), et supérieure sur les titres que le
desk négocie (8.987 contre 8.962)~: la meilleure prévision selon la mesure usuelle est la moins bonne là où elle est utilisée.

\section{Pertes alignées sur le résultat}

L'alignement le plus simple consiste à ajuster la prévision sur les données qu'utilise la décision~: les moindres carrés sur les titres négociables retrouvent la
seconde caractéristique (0.83 contre une vraie valeur de 0.80) et ignore la première. Plus généralement, une perte pondérée peut donner à chaque titre un poids égal à
la taille de la position que prendrait l'optimiseur, ou par sa liquidité, et une perte sur le rendement du portefeuille plutôt que sur la
prévision les aligne parfaitement tous les deux.

\begin{definition}[Apprentissage orienté décision, couche d'optimisation différentiable]\label{def:ml:from-prediction-to-portfolio:dfl}
L'\emph{apprentissage orienté décision}\index{apprentissage orienté décision} entraîne une prévision selon la qualité des décisions prises avec
elle~: la prévision est transmise à l'optimiseur et la perte est calculée sur le résultat de la décision (Donti, Amos et Kolter,
2017). Elle requiert une \emph{couche d'optimisation différentiable}\index{couche d'optimisation différentiable}~: un optimiseur dont la solution
peut être différenciée par rapport à ses entrées, sous forme explicite comme ici, par différentiation implicite des conditions
d'optimalité d'un programme quadratique (Amos et Kolter, 2017), ou pour les problèmes convexes généraux (Agrawal et coauteurs, 2019).
\end{definition}

\section{Politiques de portefeuille paramétriques}

\begin{definition}[Politique de portefeuille paramétrique]\label{def:ml:from-prediction-to-portfolio:ppp}
Une \emph{politique de portefeuille paramétrique}\index{politique de portefeuille paramétrique} écrit directement les poids comme fonction des
caractéristiques des titres, $w_{i,t} = \theta^\top x_{i,t}/N$ autour d'un indice de référence, et choisit $\theta$ de façon à maximiser
l'utilité moyenne sur l'échantillon (Brandt, Santa-Clara et Valkanov, 2009)~; il n'y a pas de prévision, et les covariances des caractéristiques avec les rendements
n'interviennent que par le résultat du portefeuille.
\end{definition}

Le chapitre entraîne les deux par montée de gradient sur le ratio de Sharpe des rendements nets, sur des fenêtres de 100 périodes
(le \cref{lst:ml:e2e:train}), à partir des coefficients des moindres carrés sur tous les titres~: la prévision par \omterm{def:ml:from-prediction-to-portfolio:dfl}{apprentissage orienté décision} passe par la
règle moyenne-variance, la \omterm{def:ml:reinforcement-learning-foundations:mdp}{politique} paramétrique sans son lissage. Le \Cref{tab:ml:e2e:main} et la \cref{fig:ml:e2e:sharpe} donnent les
résultats du test.

\begin{table}[htbp]
\centering\small
\begin{tabular}{@{}p{5.6cm}cccc@{}}
\toprule
coefficients choisis par & Sharpe net & (é.-t., 5 paires) & Sharpe brut & rotation\\
\midrule
moindres carrés, tous les titres (\omterm{def:ml:from-prediction-to-portfolio:pto}{prédire puis optimiser}) & 0.19 & 0.09 & 0.60 & 2.95\\
moindres carrés, titres négociables & 1.52 & 0.23 & 1.92 & 1.59\\
\omterm{def:ml:from-prediction-to-portfolio:dfl}{apprentissage orienté décision} & 1.51 & 0.22 & 1.92 & 2.75\\
\omterm{def:ml:from-prediction-to-portfolio:ppp}{politique de portefeuille paramétrique} & 1.49 & 0.23 & 1.92 & 2.90\\
coefficients vrais & 1.57 & 0.24 & 1.97 & 1.50\\
\bottomrule
\end{tabular}
\caption{Ratios de Sharpe annualisés sur les panels de test (52 périodes par an) et rotation moyenne par période (somme des variations absolues
de poids), moyennes sur cinq paires apprentissage-test. Données~: \texttt{ml\_e2e.results}.}\label{tab:ml:e2e:main}
\end{table}

\section{Apprendre à travers l'optimiseur}

Apprendre à travers la décision récupère ce que \omterm{def:ml:from-prediction-to-portfolio:pto}{prédire puis optimiser} perd~: 1.51 contre 0.19, soit un écart de 1.33 avec un écart-type
de 0.18 sur les cinq paires. Il apprend la bonne direction, 1.43 sur la seconde caractéristique et presque rien sur la
première~; l'échelle est plus grande que pour la vraie valeur car un ratio de Sharpe n'en dépend pas, et la rotation est plus élevée pour la
même raison. Mais elle ne fait pas mieux que les moindres carrés sur les titres négociables, qui sont plus simples, plus rapides et font moins de rotation.
Le gain vient de l'alignement des données sur la décision, et non de la différentiation à travers l'optimiseur~; ici l'optimiseur
est assez simple pour que les deux se confondent. L'\omterm{def:ml:from-prediction-to-portfolio:dfl}{apprentissage orienté décision} ne justifie sa complexité que lorsque la sensibilité de la décision à la
prévision varie d'une façon qu'un poids ne peut exprimer~: des contraintes qui saturent pour certains titres et pas pour d'autres, des coûts qui dépendent de la
transaction, un risque qui couple les titres.

\begin{omfigure}
\begin{tikzpicture}
\begin{axis}[width=12cm,height=6cm,ybar,bar width=10pt,ymin=0,ymax=2.1,ylabel={ratio de Sharpe net (test)},area legend,
  xtick={0,1,2,3,4},xticklabels={{prédire puis\\optimiser},{moindres\\carrés\\négociables},{apprentissage\\orienté\\décision},{politique\\paramétrique},{vrai}},xticklabel style={font=\footnotesize},
  xticklabel style={align=center},
  tick label style={font=\small},label style={font=\small},xmin=-0.5,xmax=4.5,grid=major,grid style={gray!25},
  legend columns=2,legend style={font=\footnotesize,at={(0.5,-0.3)},anchor=north,draw=none,fill=none,/tikz/every even column/.append style={column sep=8pt}},table/col sep=comma]
\addplot[fill=omDef,draw=none,bar shift=-5.5pt,error bars/.cd,y dir=both,y explicit]
  table[x=i,y=net,y error=sd]{figdata/ml/22-from-prediction-to-portfolio/sharpe.csv};
\addplot[fill=omThm,draw=none,bar shift=5.5pt,error bars/.cd,y dir=both,y explicit]
  table[x=i,y=netover,y error=sdover]{figdata/ml/22-from-prediction-to-portfolio/sharpe.csv};
\legend{{1\,500 périodes, 4 caractéristiques},{300 périodes, 20 caractéristiques}}
\end{axis}
\end{tikzpicture}

\omcaption{Ratio de Sharpe net sur les panels de test selon la manière dont les coefficients de la prévision ont été choisis, moyenne et écart-type sur
cinq couples apprentissage-test, avec des données abondantes et avec peu de données et de nombreuses variables inutiles. Données~:
\texttt{ml\_e2e.summary}.}\label{fig:ml:e2e:sharpe}
\end{omfigure}

\begin{omfigure}
\begin{tikzpicture}
\begin{axis}[width=11cm,height=5.6cm,ybar,bar width=10pt,ylabel={coefficient (en unités de 0.1~\%)},area legend,
  xtick={0,1,2,3,4},xticklabels={{prédire puis\\optimiser},{moindres\\carrés\\négociables},{apprentissage\\orienté\\décision},{politique\\paramétrique},{vrai}},xticklabel style={font=\footnotesize},
  xticklabel style={align=center},
  tick label style={font=\small},label style={font=\small},xmin=-0.5,xmax=4.5,grid=major,grid style={gray!25},
  legend columns=2,legend style={font=\footnotesize,at={(0.5,-0.3)},anchor=north,draw=none,fill=none,/tikz/every even column/.append style={column sep=8pt}},table/col sep=comma]
\addplot[fill=gray,draw=none,bar shift=-5.5pt] table[x=i,y=b1]{figdata/ml/22-from-prediction-to-portfolio/coefs.csv};
\addplot[fill=omDef,draw=none,bar shift=5.5pt] table[x=i,y=b2]{figdata/ml/22-from-prediction-to-portfolio/coefs.csv};
\legend{première variable (signal non négociable),seconde variable (signal négociable)}
\end{axis}
\end{tikzpicture}

\omcaption{Les coefficients de la prévision sur les deux variables informatives, moyennes sur cinq panels d'apprentissage. Données~:
\texttt{ml\_e2e.results}.}\label{fig:ml:e2e:coefs}
\end{omfigure}

\section{Quand l'apprentissage de bout en bout souffre de surapprentissage}

Une perte calculée sur le ratio de Sharpe net du portefeuille est un résumé bruité et de faible dimension des données~; ajuster beaucoup de paramètres à
un tel résumé revient à ajuster du bruit. Avec 300 périodes d'apprentissage et seize variables inutiles ajoutées, les moindres carrés sur les noms négociables
atteignent encore 1.04 hors échantillon, l'\omterm{def:ml:from-prediction-to-portfolio:dfl}{apprentissage orienté décision} 0.89 et la \omterm{def:ml:reinforcement-learning-foundations:mdp}{politique} paramétrique 0.87 (\cref{fig:ml:e2e:sharpe}),
avec un turnover plus élevé (3.49 contre 1.88)~: les ajustements de bout en bout ont dépensé leur latitude sur des variables qui avaient rapporté par hasard dans
l'échantillon d'apprentissage. Même avec 1\,500 périodes, les coefficients de l'\omterm{def:ml:from-prediction-to-portfolio:dfl}{apprentissage orienté décision} sur les deux variables inutiles du
premier panel valaient $-0.44$ et 0.20, contre $-0.22$ et 0.13 pour les moindres carrés sur les noms négociables. Les remèdes sont ceux de tout ajustement
sur petit échantillon~: moins de paramètres, des pénalités, un \omterm{def:ml:trees-and-boosting:early}{arrêt précoce} sur un bloc de validation et un point de départ en moindres carrés que
l'ajustement de bout en bout doit battre sur des données hors échantillon.

\begin{method}[De la prévision au portefeuille]\label{met:ml:e2e:method}
\begin{enumerate}
\item Ajuster les prévisions sur les noms, les périodes et les horizons qu'utilise le portefeuille, pondérées selon le degré de dépendance du portefeuille à leur égard.
\item Évaluer les prévisions par la performance nette du portefeuille construit à partir d'elles, jamais par la seule erreur de prévision.
\item Recourir à l'\omterm{def:ml:from-prediction-to-portfolio:dfl}{apprentissage orienté décision} ou aux \omterm{def:ml:reinforcement-learning-foundations:mdp}{politiques} paramétriques lorsque la sensibilité de la décision à la prévision ne peut pas s'exprimer
sous forme de poids, avec une forte régularisation et un indice de référence en moindres carrés.
\item Rapporter la dispersion sur des échantillons indépendants~; un écart inférieur à celle-ci n'en est pas un.
\end{enumerate}
\end{method}

\section{Tutoriel~: la perte qui rapporte}\label{tut:ml:from-prediction-to-portfolio:tut}

\begin{tutorial}
\textbf{Objectif.} Comparer \omterm{def:ml:from-prediction-to-portfolio:pto}{prédire puis optimiser}, moindres carrés alignés, \omterm{def:ml:from-prediction-to-portfolio:dfl}{apprentissage orienté décision} et une \omterm{def:ml:reinforcement-learning-foundations:mdp}{politique} paramétrique sur cinq
couples apprentissage-test, puis avec peu de données et de nombreuses variables. \textbf{État final~:} le \cref{tab:ml:e2e:main}, les \cref{fig:ml:e2e:sharpe,fig:ml:e2e:coefs}.
Les \cref{fig:ml:e2e:sharpe,fig:ml:e2e:coefs}.
\begin{tutsteps}
\item \textbf{La règle de portefeuille et son évaluation.}
\omcode{code/firm/e2eport/firm_e2eport.py}{45}{68}{Une règle moyenne-variance différentiable avec un proxy de coûts, et l'évaluation du rendement net.}\label{lst:ml:e2e:run}
\item \textbf{Apprendre à travers la décision.}
\omcode{code/firm/e2eport/firm_e2eport.py}{76}{100}{Montée de gradient sur le ratio de Sharpe net à travers la règle.}\label{lst:ml:e2e:train}
\item \textbf{Lancer} \texttt{ml\_e2e.results()}, \texttt{overfit()}, \texttt{summary()}, \texttt{forecast\_errors()} et
\texttt{fig\_e2e.py} (environ 45 secondes sur un cœur).
\end{tutsteps}
\textbf{Pour aller plus loin.} Ajouter une limite de position par nom qui devient contraignante pour certains noms et vérifier si les moindres carrés alignés
peuvent encore égaler l'\omterm{def:ml:from-prediction-to-portfolio:dfl}{apprentissage orienté décision}~; ajouter un \omterm{def:ml:trees-and-boosting:early}{arrêt précoce} sur un bloc de validation à l'ajustement de bout en bout.
\end{tutorial}

\section{Construction~: portefeuilles de bout en bout}\label{bld:ml:from-prediction-to-portfolio:e2eport}

\begin{build}
\bfield{Rôle} Prévisions jugées et, lorsque cela rapporte, entraînées par les portefeuilles qu'elles produisent.
\bfield{Interface} \texttt{char\_panel(seed, T, N, K, n\_noise, phi)}, \texttt{mv\_layer(mu, w\_prev, kappa, gamma, sigma)},
\texttt{run(b, X, R, tradeable, cost, gamma, sigma, policy)}, \texttt{least\_squares(X, R, rows)},
\texttt{train\_decision(X, R, tradeable, init, policy, epochs, lr, window, seed)}.
\bfield{Règles} Chaque résultat de bout en bout est rapporté par rapport aux moindres carrés alignés sur les mêmes données hors échantillon, avec sa dispersion sur
échantillons indépendants.
\bfield{Tests d'acceptation} \texttt{code/firm/e2eport/tests/}~: la couche en forme close est égale à la solution \texttt{firm.portcons}
du livre~7 lorsque les coûts sont nuls~; le gradient de la couche coïncide avec les différences finies~; l'évaluation impute les coûts aux variations de poids~;
l'\omterm{def:ml:from-prediction-to-portfolio:dfl}{apprentissage orienté décision} sur un problème sans bruit retrouve la direction des vrais coefficients.
\bfield{Pour aller plus loin} Une couche générale de programme quadratique par différentiation implicite, avec un modèle de risque factoriel~; limites par nom~;
\omterm{def:ml:trees-and-boosting:early}{arrêt précoce}.
\end{build}

\begin{omsources}
\item M.~W.~Brandt, P.~Santa-Clara et R.~Valkanov, «~Parametric portfolio policies~: exploiting characteristics in the
cross-section of equity returns~», \emph{Review of Financial Studies} 22(9), 2009.
\item A.~N.~Elmachtoub et P.~Grigas, «~Smart `predict, then optimize'~», \emph{Management Science} 68(1), 2022.
\item P.~L.~Donti, B.~Amos et J.~Z.~Kolter, «~Task-based end-to-end model learning in stochastic optimization~»,
arXiv:1703.04529, 2017.
\item B.~Amos et J.~Z.~Kolter, «~OptNet: differentiable optimization as a layer in neural networks~», arXiv:1703.00443, 2017.
\item A.~Agrawal et coauteurs, «~Differentiable convex optimization layers~», arXiv:1910.12430, 2019.
\item A.~Butler et R.~H.~Kwon, «~Integrating prediction in mean-variance portfolio optimization~», \emph{Quantitative Finance}
23(3), 2023.
\end{omsources}

\section{Exercices}

\begin{exercise}[$\star$]\label{exo:ml:from-prediction-to-portfolio:1}
Dériver la règle $w = (\mu + 2\kappa w_{\mathrm{prev}})/(\gamma\sigma^2 + 2\kappa)$ à partir de l'objectif à un seul nom, puis déterminer
ce qu'elle devient avec $\kappa = 0$ et lorsque $\kappa\to\infty$.
\end{exercise}

\begin{exercise}[$\star$]\label{exo:ml:from-prediction-to-portfolio:2}
Déterminer, avec les paramètres du chapitre, quelle fraction de l'écart à sa cible le poids d'un nom comble à chaque période.
\end{exercise}

\begin{exercise}[$\star$]\label{exo:ml:from-prediction-to-portfolio:3}
Expliquer pourquoi une perte fondée sur le ratio de Sharpe laisse indéterminée l'échelle des coefficients, et quel effet cela a sur le turnover.
\end{exercise}

\begin{exercise}[$\star\star$]\label{exo:ml:from-prediction-to-portfolio:4}
Expliquer pourquoi la prévision par moindres carrés sur l'ensemble des noms a une erreur globale plus faible et une erreur plus élevée sur les noms négociables.
\end{exercise}

\begin{exercise}[$\star\star$]\label{exo:ml:from-prediction-to-portfolio:5}
Déterminer dans quels cas l'\omterm{def:ml:from-prediction-to-portfolio:dfl}{apprentissage orienté décision} battrait les moindres carrés sur les noms négociables dans ce cadre.
\end{exercise}

\begin{exercise}[$\star\star$]\label{exo:ml:from-prediction-to-portfolio:6}
\emph{Trouver la faille.} «~On a entraîné le réseau de bout en bout sur le ratio de Sharpe du portefeuille et obtenu 3.1 au backtest, contre
1.2 pour la prévision entraînée sur les rendements, donc l'entraînement de bout en bout en vaut la peine.~»
\end{exercise}

\begin{exercise}[$\star\star\star$]\label{exo:ml:from-prediction-to-portfolio:7}
\emph{Programmation.} Ajouter une pénalité ridge à l'entraînement orienté décision dans le scénario à petit échantillon et choisir son intensité sur un
bloc de validation. Réduit-elle l'écart avec les moindres carrés alignés~?
\end{exercise}

\begin{exercise}[$\star\star\star$]\label{exo:ml:from-prediction-to-portfolio:8}
Pour le problème de moyenne-variance sans contrainte $\max_w\mu^\top w - \frac\gamma2w^\top\Sigma w$, calculer $\partial w^*/\partial\mu$
et expliquer comment le gradient d'une perte de portefeuille remonte jusqu'à la prévision.
\end{exercise}

\section{Problème~: La perte qui rapporte}

\begin{problem}[{Problème du week-end --- entraîner sur ce que l'on négocie}]\label{pb:ml:from-prediction-to-portfolio:1}
Le panel du chapitre, la règle de portefeuille et les cinq façons de choisir la prévision.

\medskip\noindent\textbf{Partie I --- Le cadre.}
\begin{enumerate}
\item Que contient chacune des deux moitiés de l'univers, et que peut représenter la prévision linéaire~?
\item Que fait la règle de portefeuille, et comment les coûts sont-ils imputés~?
\item Quels coefficients les moindres carrés sur l'ensemble des noms apprennent-ils, et pourquoi~?
\item Comment les deux prévisions par moindres carrés se comparent-elles selon l'erreur quadratique moyenne~?
\end{enumerate}

\medskip\noindent\textbf{Partie II --- La comparaison.}
\begin{enumerate}[resume]
\item Donner les ratios de Sharpe nets et bruts ainsi que le turnover des cinq méthodes.
\item Quel est l'écart entre l'\omterm{def:ml:from-prediction-to-portfolio:dfl}{apprentissage orienté décision} et \omterm{def:ml:from-prediction-to-portfolio:pto}{prédire puis optimiser}, et quelle est sa dispersion~?
\item Pourquoi les moindres carrés alignés et l'\omterm{def:ml:from-prediction-to-portfolio:dfl}{apprentissage orienté décision} font-ils jeu égal~?
\item Pourquoi le turnover de la \omterm{def:ml:reinforcement-learning-foundations:mdp}{politique} paramétrique est-il le plus élevé~?
\end{enumerate}

\medskip\noindent\textbf{Partie III --- \omterm{def:ml:why-financial-machine-learning-is-different:generalisation}{Surapprentissage}.}
\begin{enumerate}[resume]
\item Que se passe-t-il avec 300 périodes et vingt variables~?
\item Quels coefficients révèlent le \omterm{def:ml:why-financial-machine-learning-is-different:generalisation}{surapprentissage}~?
\item Quelle régularisation utiliserait-on~?
\item Comment choisir entre les méthodes sur des données réelles~?
\end{enumerate}

\medskip\noindent\textbf{Partie IV --- Le verdict.}
\begin{enumerate}[resume]
\item Énoncer le \emph{résultat nommé}~: le ratio de Sharpe net de \omterm{def:ml:from-prediction-to-portfolio:pto}{prédire puis optimiser} face à l'\omterm{def:ml:from-prediction-to-portfolio:dfl}{apprentissage orienté décision} et à la
\omterm{def:ml:from-prediction-to-portfolio:ppp}{politique de portefeuille paramétrique}, avec la dispersion de l'écart sur les graines aléatoires.
\item Quand \omterm{def:ml:from-prediction-to-portfolio:pto}{prédire puis optimiser} suffit-il~?
\item Qu'est-ce qui rendrait l'\omterm{def:ml:from-prediction-to-portfolio:dfl}{apprentissage orienté décision} à la hauteur de son coût~?
\item Quel lien cela entretient-il avec le coefficient de transfert~?
\item Que doit contenir un rapport de recherche sur une nouvelle prévision, après ce chapitre~?
\item Où les coûts entrent-ils dans la perte d'entraînement, et où devraient-ils entrer~?
\item Comment procéder avec une prévision par \omterm{def:ml:neural-networks-for-noisy-tabular-data:mlp}{réseau de neurones}~?
\item En une phrase~: pour quoi doit-on entraîner une prévision~?
\end{enumerate}
\end{problem}

\section{Questions d'entretien}

\begin{interviewq}[$\star$ \iqroles{chercheur}]\label{iq:ml:from-prediction-to-portfolio:1}
Pourquoi une prévision plus précise peut-elle produire un portefeuille moins bon~?
\end{interviewq}

\begin{interviewq}[$\star\star$ \iqroles{chercheur, ingénieur ML}]\label{iq:ml:from-prediction-to-portfolio:2}
Qu'est-ce que l'\omterm{def:ml:from-prediction-to-portfolio:dfl}{apprentissage orienté décision}, et comment différenciez-vous à travers un optimiseur~?
\end{interviewq}

\begin{interviewq}[$\star\star$ \iqroles{chercheur}]\label{iq:ml:from-prediction-to-portfolio:3}
Décrivez les \omterm{def:ml:from-prediction-to-portfolio:ppp}{politiques de portefeuille paramétriques} de Brandt, Santa-Clara et Valkanov. Quels sont leurs avantages et leurs risques~?
\end{interviewq}

\begin{interviewq}[$\star\star$ \iqroles{chercheur, trader}]\label{iq:ml:from-prediction-to-portfolio:4}
Comment intégreriez-vous les coûts de transaction dans l'entraînement d'un modèle alpha~?
\end{interviewq}

\begin{interviewq}[$\star\star$ \iqroles{ingénieur ML}]\label{iq:ml:from-prediction-to-portfolio:5}
Un réseau de portefeuille de bout en bout obtient au backtest des résultats bien meilleurs que \omterm{def:ml:from-prediction-to-portfolio:pto}{prédire puis optimiser}. Que vérifiez-vous~?
\end{interviewq}

\begin{interviewq}[$\star\star\star$ \iqroles{chercheur}]\label{iq:ml:from-prediction-to-portfolio:6}
Montrez comment les conditions KKT d'un programme quadratique donnent la dérivée de sa solution par rapport au terme linéaire.
\end{interviewq}
